% Options for packages loaded elsewhere
\PassOptionsToPackage{unicode}{hyperref}
\PassOptionsToPackage{hyphens}{url}
\PassOptionsToPackage{dvipsnames,svgnames,x11names}{xcolor}
\documentclass[
  11pt,
  letterpaper,
]{article}
\usepackage{xcolor}
\usepackage[margin=0.5in]{geometry}
\usepackage{amsmath,amssymb}
\setcounter{secnumdepth}{-\maxdimen} % remove section numbering
\usepackage{iftex}
\ifPDFTeX
  \usepackage[T1]{fontenc}
  \usepackage[utf8]{inputenc}
  \usepackage{textcomp} % provide euro and other symbols
\else % if luatex or xetex
  \usepackage{unicode-math} % this also loads fontspec
  \defaultfontfeatures{Scale=MatchLowercase}
  \defaultfontfeatures[\rmfamily]{Ligatures=TeX,Scale=1}
\fi
\usepackage{lmodern}
\ifPDFTeX\else
  % xetex/luatex font selection
  \setmainfont[]{Arial}
\fi
% Use upquote if available, for straight quotes in verbatim environments
\IfFileExists{upquote.sty}{\usepackage{upquote}}{}
\IfFileExists{microtype.sty}{% use microtype if available
  \usepackage[]{microtype}
  \UseMicrotypeSet[protrusion]{basicmath} % disable protrusion for tt fonts
}{}
\makeatletter
\@ifundefined{KOMAClassName}{% if non-KOMA class
  \IfFileExists{parskip.sty}{%
    \usepackage{parskip}
  }{% else
    \setlength{\parindent}{0pt}
    \setlength{\parskip}{6pt plus 2pt minus 1pt}}
}{% if KOMA class
  \KOMAoptions{parskip=half}}
\makeatother
\usepackage{longtable,booktabs,array}
\newcounter{none} % for unnumbered tables
\usepackage{calc} % for calculating minipage widths
% Correct order of tables after \paragraph or \subparagraph
\usepackage{etoolbox}
\makeatletter
\patchcmd\longtable{\par}{\if@noskipsec\mbox{}\fi\par}{}{}
\makeatother
% Allow footnotes in longtable head/foot
\IfFileExists{footnotehyper.sty}{\usepackage{footnotehyper}}{\usepackage{footnote}}
\makesavenoteenv{longtable}
\setlength{\emergencystretch}{3em} % prevent overfull lines
\providecommand{\tightlist}{%
  \setlength{\itemsep}{0pt}\setlength{\parskip}{0pt}}
\usepackage{bookmark}
\IfFileExists{xurl.sty}{\usepackage{xurl}}{} % add URL line breaks if available
\urlstyle{same}
\hypersetup{
  colorlinks=true,
  linkcolor={black},
  filecolor={black},
  citecolor={black},
  urlcolor={black},
  pdfcreator={LaTeX}}

\author{}
\date{}

\begin{document}


\section{Response to Reviewers, Round
2}\label{response-to-reviewers-round-2}

\textbf{Manuscript:} Validated Synthetic Patient Generation for Small
Longitudinal Cohorts: Coagulation Dynamics Across Pregnancy\\
\textbf{Journal:} npj Systems Biology and Applications\\
\textbf{Decision:} Minor revision\\
\textbf{Date:} 16 November 2026

\begin{center}\rule{0.5\linewidth}{0.5pt}\end{center}

\subsection{Note on this document}\label{note-on-this-document}

Each reviewer comment is quoted below. Each comment is followed by our
response and the location of the corresponding manuscript change. Text
added in this round appears in green in the marked manuscript
(\texttt{\textbackslash{}rthree\{...\}}). Text added in the first round
now appears in black. The supplement is supplied without color, and each
response lists its new supplement text and tables. A
\texttt{\textbackslash{}reviewmode} setting in the preamble changes all
colored text to black for the final version.

We also corrected the Data Availability statement. The de-identified
real patient data are in the public project repository, and the
statement now says so.

\begin{center}\rule{0.5\linewidth}{0.5pt}\end{center}

\section{Reviewer 1}\label{reviewer-1}

\begin{quote}
The reviewer thanks the authors for the effort in making the revisions.
The authors have answered all my questions and the manuscript is now
suitable for publication.
\end{quote}

We thank Reviewer 1 for the careful reading in both rounds. The comments
improved the privacy analysis, the robustness tests, and the description
of the method.

\begin{center}\rule{0.5\linewidth}{0.5pt}\end{center}

\section{Reviewer 3}\label{reviewer-3}

We thank Reviewer 3 for the careful reading and for the three remaining
comments. We added new analyses for comments 1 and 3 and report them in
a new Supplement subsection, ``Supplementary Bandwidth and Positivity
Results.''

\subsubsection{R3.1: Advantage over a Gaussian mixture, and the relation
to kernel-density
estimation}\label{r3.1-advantage-over-a-gaussian-mixture-and-the-relation-to-kernel-density-estimation}

\begin{quote}
The revised results show that simpler approaches, particularly the
Gaussian mixture and copula models, achieve comparable or better
performance on several metrics. The authors should clarify the specific
methodological advantage of stochastic attention, particularly given its
exact representation as a finite Gaussian mixture. The relationship
between stochastic attention and conventional Gaussian kernel-density
estimation deserves further discussion.
\end{quote}

\emph{Stochastic attention did not perform better than the fitted
Gaussian mixture or the copula on these metrics, and we do not claim
that it did.} Its advantage is in how the distribution is built and how
it can be extended. The mixture has one component centered on each
stored profile, every component has the same isotropic variance, and the
inverse temperature is the only parameter. We chose the inverse
temperature from the attention-entropy transition, without held-out
data. The fitted two-component mixture instead estimated means and full
covariance matrices in \(18\) dimensions from \(23\) profiles.

The relation to kernel-density estimation is exact. With equal
multiplicity weights, the stochastic-attention distribution is an
isotropic Gaussian kernel-density estimate on the \(23\) unit-normalized
profiles, with kernel width \(h=\beta^{-1/2}\). With unequal weights, it
is a weighted kernel-density estimate, which is the form we used for
conditional generation. Stochastic attention differs from conventional
kernel-density estimation in two ways. First, it places the kernels on
unit-normalized profiles and draws each profile's magnitude separately.
Second, it sets the kernel width from the entropy transition instead of
a bandwidth rule.

To test the second difference, we compared the entropy transition with
three standard bandwidth rules on the same profiles. The entropy
transition gave \(h=0.58\). Leave-one-out likelihood gave \(0.29\),
Scott's rule gave \(0.21\), and Silverman's rule gave \(0.19\). These
narrower kernels corresponded to \(4.1\)--\(9.0\) times \(\beta^*\).
Profiles drawn with the narrower kernels stayed closer to the stored
profiles. Median novelty fell from \(0.50\) to \(0.21\)--\(0.34\), and
the median distance to the closest real record fell from \(13.9\) to
\(9.9\)--\(11.9\). Cross-visit correlation error fell from \(0.60\) to
\(0.43\)--\(0.50\), and median relative error rose from \(1.7\%\) to
\(2.1\)--\(2.2\%\). Mechanistic overlap was \(0.91\)--\(0.92\) for every
width. Of the four widths, the entropy transition gave the widest
kernel, the most novel profiles, and the lowest marginal error.

The Gaussian-mixture identity also explains how the method extends.
Stochastic attention defines an energy on the stored profiles, and the
Langevin sampler acts on that energy. With no other terms in the energy,
the target is the Gaussian mixture. We used this case in the first round
to check the Langevin sampler against exact draws. A constraint term
\(\lambda C\) can be added to the energy, and the same sampler applies.
A quadratic term keeps the target a Gaussian mixture, but a
non-quadratic term does not. We tested a non-quadratic positivity term
(R3.3). With this term, Langevin sampling agreed with exact rejection
sampling from the same target at every constraint weight.

\textbf{Changes:} Methods, ``Multiplicity-Weighted Stochastic
Attention'' (one sentence stating the kernel-density form and kernel
width); Methods, ``Baseline and Validation Framework'' (new paragraph
describing the bandwidth comparison, the positivity term, and the log
representations); Results, ``Marginal Plausibility'' (how the
kernel-density baseline differed from the kernel-density form of
stochastic attention); Discussion (three sentences added to the
paragraph on the PCA representation and the generator); Supplement, new
subsection ``Supplementary Bandwidth and Positivity Results'' and
Supplementary Table 17. Green. \texttt{\textbackslash{}rthree\{...\}}

\begin{center}\rule{0.5\linewidth}{0.5pt}\end{center}

\subsubsection{R3.2: Membership-inference AUC and sharing synthetic
data}\label{r3.2-membership-inference-auc-and-sharing-synthetic-data}

\begin{quote}
The membership-inference analysis reports an AUC of approximately 0.97,
indicating substantial privacy risk. The authors should clarify the
implications for sharing synthetic patient data and distinguish
exploratory computational applications from privacy-preserving release.
\end{quote}

\emph{We agree. An AUC of \(0.97\) means that the synthetic cohort does
not protect membership, so stochastic attention as used here is not a
privacy-preserving release method.} Someone who already holds a
candidate de-identified profile can usually tell whether that profile
was stored in the memory matrix. The revised Discussion now states this
directly. The synthetic cohorts support exploratory computational work,
such as hypothesis generation, mechanistic modeling, and workflow
testing. A cohort generated from restricted data should be shared under
the same terms as those data. Releasing synthetic data from sensitive
cohorts would require a formal privacy method, such as differential
privacy, which we did not test.

The membership signal came from the target distribution, not from the
numerical sampler. Direct samples from the exact mixture gave a mean AUC
of \(0.98\). In the first round, lowering the inverse temperature moved
synthetic profiles farther from the stored profiles but reduced
cross-visit covariance fidelity. The new bandwidth comparison (R3.1)
agreed. The standard bandwidth rules chose narrower kernels, which
placed profiles closer to the stored profiles, with median distances to
the closest real record of \(9.9\)--\(11.9\) compared with \(13.9\) at
\(\beta^*\). No kernel width in this comparison gave a formal privacy
guarantee.

In this study, the de-identified source data are public in the project
repository. The synthetic cohorts therefore expose no profile that is
not already public.

\textbf{Changes:} Discussion (four sentences added to the limitations
paragraph, after the bounded statement of what the membership test
showed); Data Availability (corrected). Green.
\texttt{\textbackslash{}rthree\{...\}}

\begin{center}\rule{0.5\linewidth}{0.5pt}\end{center}

\subsubsection{R3.3: Biological plausibility and
positivity}\label{r3.3-biological-plausibility-and-positivity}

\begin{quote}
Only 54\% of the generated profiles satisfy the biological plausibility
criteria, and 24\% satisfy both biological and strict mechanistic
criteria. The authors should emphasize this limitation and discuss
whether positivity-preserving transformations or constrained
reconstruction are necessary for future applications.
\end{quote}

\emph{We agree that this is a limitation, and the revised Discussion now
states it with the numbers. For future applications that use hormone
values, a positivity-preserving transform is needed, and the log
transform is the one we recommend.} We tested both remedies the reviewer
named.

Every profile that failed the biological check had a negative estradiol
or progesterone value. Linear decoding has unbounded support, and the
hormones lie closer to zero relative to their variation than the
coagulation factors do. The strict mechanistic check is also demanding
for real patients. It requires all \(15\) predicted-to-recorded TGA
ratios (five features at three visits) to fall within the real 5th to
95th percentile bands. When each real profile was checked against bands
built from the other \(22\) real profiles, only \(7\) of \(23\) passed.
With bands built from all \(23\) real profiles, which were the bands
applied to the synthetic profiles, \(8\) of \(23\) passed. The \(24\%\)
of synthetic profiles that passed both checks should be read against
this real-patient rate of \(30\)--\(35\%\).

We first tested constrained reconstruction. We added a positivity term
\(\lambda C\) to the energy. The term \(C\) sums the squared negative
parts of the six decoded hormone values, each divided by its real-cohort
standard deviation. We sampled with Langevin dynamics at \(\lambda=0\),
\(1\), \(10\), and \(100\), and compared each run with exact rejection
sampling from the same target. Langevin and exact samples agreed at
every \(\lambda\): two-sample Kolmogorov-Smirnov tests on \(1{,}000\)
draws per sampler gave \(p\ge0.002\) for all \(24\) comparisons, above
the Bonferroni threshold of \(0.0004\). As \(\lambda\) increased, the
number of profiles with a negative hormone value fell from about \(50\)
at \(\lambda=0\) to \(4\)--\(6\) at \(\lambda=100\). A hard constraint
removed every negative value and kept \(51\%\) of exact draws. The
constraint also distorted the hormone distributions. Removing the
negative tail raised the hormone means, so the median hormone error rose
from about \(5\%\) to \(30\%\), and the median relative error across all
features rose from \(1.2\)--\(1.3\%\) to \(2.5\%\).

We then tested a positivity-preserving transform. We replaced the six
hormone values with their logarithms before standardization and PCA, and
exponentiated them after decoding. All \(100\) profiles passed the
biological check, and \(45\) also passed the strict mechanistic check.
The untransformed representation, sampled the same way, gave \(48\) and
\(20\). Median relative error was \(1.4\%\), compared with \(1.7\%\)
without the transform. Median hormone error was \(6.2\%\), compared with
\(6.4\%\). Cross-visit correlation error was \(0.58\), compared with
\(0.60\), and mechanistic overlap was \(0.90\), compared with \(0.91\).
Log-transforming all \(72\) assays also removed the negative values and
gave a median relative error of \(1.3\%\), but mechanistic overlap fell
to \(0.88\).

The log transform of the hormones therefore removed the negative values
without a loss in fidelity, while the constraint term removed them at a
cost in hormone accuracy. The reported cohort and the results in the
main text are unchanged. The log representation is reported in the
Supplement and recommended in the Discussion for future applications.

\textbf{Changes:} Results, ``Interpolation and Extrapolation'' (replaced
the first-round sentence on positive-variable transformations with the
real-patient pass rate and the log-transform result); Discussion (new
sentences in the paragraph on the linear representation, stating the
\(54/100\) and \(24/100\) rates, the real-patient rate, both remedies,
and the recommendation); Supplementary Table 14 (new column for the real
patients); Supplementary Table 18 and the new subsection ``Supplementary
Bandwidth and Positivity Results.'' Green.
\texttt{\textbackslash{}rthree\{...\}}

\end{document}
